\scanpage{42}
mais il est facile de le retrouver, en choisissant un gén $u$ de $\mu$, et en introduisant les valeurs propres $\zeta,\zeta^{-1}$ de $u$ (satisfaisant $\zeta^p=1$ donc $\zeta\ne\zeta^{-1}$ car $p$ impair, $\zeta\ne\pm1$ \inserted{\unsure{il suffirait de|il suffisant de|racines de} $p\unsure{>|=|\geq}2$, \unsure{premier pourr.|premiè. pourr.|primitives}} puisque repr. \uline{fid} de $\mathbb Z/p\mathbb Z$) le commutant \inserted{de $u$} est formé des $(\lambda,\lambda^{-1})$, le normalisateur est formé des matrices qui sont soit de cette forme diagonale, soit le produit d'une telle par une symétrie \ldots\ Ici $N/H\simeq\{\pm1\}_S\subset(\mathbb Z/p\mathbb Z)^\times_{\deleted{\unsure{S}}S}$

Si $p=\mathrm{car}\,k$, alors $u$ est de la forme
\[
 u=\mathrm{id}_V+v
\]
avec \deleted{\unsure{$v^2$.}}
\[
 v^2=0,\quad v\ne0
\]
et le commutant est formé des matrices \uline{unimodulaires} de la forme \deleted{\unsure{$\lambda$}} $\alpha\,\mathrm{id}+\beta v$ \deleted{il faut $\alpha$ inv. et $\unsure{v|\alpha}^2=1$} \deleted{\unsure{donc le}} la condition unim. signifie $\alpha^2=1$ (condition étale \inserted{$\alpha=\pm1$}, puisque $p\ne2$ !), et on a encore la lissité, la connexité des $U$, etc. \deleted{\unsure{Si c'est ceci} fait que de}

b) \uline{Cas général} Il suffit de prouver que $N$ est lisse et $N^0$ commutatif -- Comme \deleted{$N/N^0$ est un grpe étale $\subset$} l'image inverse par $N\longrightarrow\underline{\operatorname{Aut}}(\mu)=(\mathbb Z/p\mathbb Z)^\times_S$ de la section unité est un sous-groupe ouvert de $N$ qui coïncide avec $N^0$ en chaque fibre, il est égal à $N^0$, de plus $N/N^0\subset(\mathbb Z/p\mathbb Z)^\times_S$ est
% archive sheet number: 41
